\documentclass[final]{beamer}
\usefonttheme[onlymath]{serif}

% \usetheme{SUPoster}
\usetheme{CCRMAPoster}


\usepackage[orientation=landscape,size=a0,scale=1.6]{beamerposter}
\usepackage[absolute,overlay]{textpos}
\setlength{\TPHorizModule}{1cm}
\setlength{\TPVertModule}{1cm}


\usepackage{lmodern} % needed to make math mode fonts scale properly
\usepackage{url}
\usepackage{graphics}
\usepackage{bold-extra}



\usepackage{enumitem}
% \setitemize[1]{label=$\blacktriangleright$}
% \setlist[itemize,1]{label=$\blacktriangleright$}

\setlist[itemize]{label=\usebeamerfont*{itemize item}%
  \usebeamercolor[fg]{itemize item}
  \usebeamertemplate{itemize item},
  leftmargin=2em}

% next line is necessary to make beamer and enumitem play nicely together (see http://tex.stackexchange.com/questions/45921/make-enumerate-have-beamer-themes-when-using-enumitem/45950#45950)
\setlist[enumerate]{leftmargin=2.2em,%
label=\protect\usebeamerfont{enumerate item}%
    \protect\usebeamercolor[fg]{enumerate item}%
    \insertenumlabel. }




%% DOCUMENT DATA

\title{The Title}
\subtitle{And the subtitle}

\author{Author 1 \and Author 2}
\institute{Center for Computer Research in Music and Acoustics - Stanford University}

\logoleft{\begin{center} \includegraphics[height=8cm]{images/SU_Seal_Blk_pos.eps}\end{center}}
\logoright{\begin{center} \includegraphics[height=2.5cm]{images/ccrma-complete-mod1.pdf}\end{center}}

\footer{The Conference Name, The Place - The Date \hfill \url{http://ccrma.stanford.edu/}}
\date{}





\begin{document}


\definecolor{almostWhite}{RGB}{240,240,240}
\definecolor{lightGray}{RGB}{220,220,220}
\definecolor{tbllinecolor}{RGB}{200,200,200}
\beamertemplateshadingbackground{lightGray}{almostWhite}
\newcommand{\vltbl}{{\color{tbllinecolor}\vrule}}


\begin{frame}[fragile]{}



%%%%%%%%%%%%%%%%%%%%%%%%%%%%%%%%%%%%%%%%%%%%%%%%%%%%%%%%%%%%%%%%%%%%
% IMPORTANT:
%   Here is where you define your poster's layout params (widths and
%   vertical positions).
%   Feel free to add/remove more, to make your layout fit your needs
%
%   All dimensions are in cm (I think).
%%%%%%%%%%%%%%%%%%%%%%%%%%%%%%%%%%%%%%%%%%%%%%%%%%%%%%%%%%%%%%%%%%%%

% centralize columns layout
\newcommand{\vstart}{11.5}
\newcommand{\vstartCols}{20.5}
\newcommand{\fullwidth}{114.5}  % this is the key to the values below
\newcommand{\lcolwidth}{30}
\newcommand{\mcolwidth}{50.5}
\newcommand{\rcolwidth}{30}
\newcommand{\firstcolpos}{2}
\newcommand{\secondcolpos}{34}
\newcommand{\thirdcolpos}{86.5}
\newcommand{\bottomblockstart}{70.3}


% this will add some padding to the blocks, to avoid text reaching
% the border (looks bad)
\newenvironment{paddedBlock}[2][0.95\linewidth]
    {\begin{block}{#2}\begin{minipage}{#1}}
    {\end{minipage}\end{block}}






%%%%%%%%%%%%%%%%%%%%%%%%%%%%%%%%%%%%%%%%%%%%%%%%%%%%%%%%%%%%%%%%%%%%
% And now ... the content
%%%%%%%%%%%%%%%%%%%%%%%%%%%%%%%%%%%%%%%%%%%%%%%%%%%%%%%%%%%%%%%%%%%%


%%%%%%%%%%%%%%%%%%%%%%%%%
%%%%%%%%%%%%%%%%%%%%%%%%%
%% TOP ROW
%%%%%%%%%%%%%%%%%%%%%%%%%
%%%%%%%%%%%%%%%%%%%%%%%%%

\begin{textblock}{\fullwidth}(2,\vstart)
\begin{paddedBlock}[\linewidth]{Wide Block}

Here you could add some nice figures or something eye-catching

\end{paddedBlock}
\end{textblock}



%%%%%%%%%%%%%%%%%%%%%%%%%
%%%%%%%%%%%%%%%%%%%%%%%%%
%% LEFT COLUMN
%%%%%%%%%%%%%%%%%%%%%%%%%
%%%%%%%%%%%%%%%%%%%%%%%%%
\begin{textblock}{\lcolwidth}(\firstcolpos,\vstartCols)

\begin{paddedBlock}{Left Column 1}

\begin{itemize}
\item A
\item Fake
\item List
\end{itemize}


\end{paddedBlock}


\begin{paddedBlock}{Left Column 2}

\begin{description}
\item[Two-step solution:]\hfill
    % \vspace{0.5cm}
    \begin{enumerate}
    \item Think about it
    \item Do it
    \end{enumerate}
\end{description}

\vspace{1cm}

\alert{Check out these equations!}

\begin{align}
1 &= \cos(2\pi) + j\sin(2\pi) \\
  &= e^{-j2\pi}
\end{align}

\end{paddedBlock}


\end{textblock}


%%%%%%%%%%%%%%%%%%%%%%%%%
%%%%%%%%%%%%%%%%%%%%%%%%%
%% MIDDLE COLUMN
%%%%%%%%%%%%%%%%%%%%%%%%%
%%%%%%%%%%%%%%%%%%%%%%%%%

\begin{textblock}{\mcolwidth}(\secondcolpos,\vstartCols)

\begin{paddedBlock}{Middle Column 1}

I run out of ideas \ldots

\end{paddedBlock}

\begin{paddedBlock}{Middle Column 2}

\ldots still waiting

\end{paddedBlock}


\end{textblock}


%%%%%%%%%%%%%%%%%%%%%%%%%
%%%%%%%%%%%%%%%%%%%%%%%%%
%% RIGHT COLUMN
%%%%%%%%%%%%%%%%%%%%%%%%%
%%%%%%%%%%%%%%%%%%%%%%%%%
\begin{textblock}{\rcolwidth}(\thirdcolpos,\vstartCols)


\begin{paddedBlock}{Right Column 1}

And that's a wrap!

\end{paddedBlock}




%%%%%%%%%%%%%%%%%%%%%%%%%
%% BIBLIOGRAPHY
%%%%%%%%%%%%%%%%%%%%%%%%%
\vspace{0.5cm}
\textbf{\normalsize References}
\bibliographystyle{plain}
% \bibliographystyle{apalike}
% \nocite{*}
\bibliography{references}

\end{textblock}



%%%%%%%%%%%%%%%%%%%%%%%%%
%%%%%%%%%%%%%%%%%%%%%%%%%
%% BOTTOM ROW
%%%%%%%%%%%%%%%%%%%%%%%%%
%%%%%%%%%%%%%%%%%%%%%%%%%

\begin{textblock}{\fullwidth}(2,\bottomblockstart)
\begin{paddedBlock}[0.98\linewidth]{Finally}

In case you need it, you can do this too

\end{paddedBlock}
\end{textblock}

\end{frame}





\end{document}